\centerline{\bf \Huge  Что? Где? Кто?}
\title{Что? Где? Кто?}

\begin{flushright}
    \scriptsize — Что вы предпочтете, Господин Супранович?\\
— Я вот пытался об этом думать, но подумал, что не буду думать.\\
Андрей Супранович. ЧГК.

\end{flushright}

\problem{1.1} В трех ящиках находятся крупа, вермишель и сахар. На первом ящике написано «крупа», на втором – «вермишель», на третьем – «крупа или сахар». Что в каком ящике находится, если содержимое каждого из ящиков не соответствует надписи на нем?

\problem{1.2}Когда три подруги — Надя, Валя и Маша — вышли гулять, на них были белое, красное и синее платья. Туфли их были тех же трех цветов, но только у Нади цвета туфель и платья совпадали. При этом у Вали ни платье, ни туфли не были синими, а Маша была в красных туфлях. Определите цвет платьев и туфель каждой из подруг.

\problem{1.3}На столе лежат в ряд четыре фигуры: треугольник, круг, прямоугольник и ромб. Они окрашены в разные цвета: красный, синий, жёлтый, зелёный. Известно, что красная фигура лежит между синей и зелёной; справа от жёлтой фигуры лежит ромб; круг лежит правее и треугольника и ромба; треугольник лежит не с краю; синяя и жёлтая фигуры лежат не рядом. Определите, в каком порядке лежат фигуры и какого они цвета.

\problem{2} Однажды Миша, Витя и Коля заметили, что принесли в детский сад одинаковые игрушечные машинки. У Миши есть машинка с прицепом, есть маленькая машинка и есть зеленая машинка без прицепа. У Вити есть машинка без прицепа и маленькая зеленая с прицепом, а у Коли — большая машинка и маленькая синяя с прицепом. Машинку какого вида (по цвету, размеру и наличию прицепа) принесли мальчики в детский сад? Ответ объясните.

\problem{3}Иван, Петр и Сидор ели конфеты. Их фамилии – Иванов, Петров и Сидоров. Иванов съел на 2 конфеты меньше Ивана, Петров – на 2 конфеты меньше Петра, а Петр съел больше всех. У кого из них какая фамилия?

\problem{4}На улице, став в кружок, разговаривают четыре девочки: Аня, Валя, Галя и Нина. Девочка в зеленом платье (не Аня и не Валя) стоит между девочкой в голубом платье и Ниной. Девочка в белом платье стоит между девочкой в розовом платье и Валей. Какое платье на каждой из девочек?

\newpage

\problem{5}В очереди в школьный буфет стоят Вика, Соня, Боря, Денис и Алла. Вика стоит впереди Сони, но после Аллы; Боря и Алла не стоят рядом; Денис не находится рядом ни с Аллой, ни с Викой, ни с Борей. В каком порядке стоят ребята?


\problem{6} В забеге шести спортсменов Андрей отстал от Бориса и еще от двух спортсменов. Виктор финишировал после Дмитрия, но ранее Геннадия. Дмитрий опередил Бориса, но все же пришел после Евгения. Какое место занял каждый спортсмен?

\problem{7}Четыре подруги пришли на каток, каждая со своим братом. Они разбились на пары и начали кататься. Оказалось, что в каждой паре «кавалер» выше «дамы» и никто не катается со своей сестрой. Самым высоким в компании был Юра Воробьев, следующим по росту — Андрей Егоров, потом Люся Егорова, Сережа Петров, Оля Петрова, Дима Крымов, Инна Крымова и Аня Воробьева. Определите, кто с кем катался?